\MathBookSourcePage{144}
\subsection{纯 Neumann 问题的变分表述}
\label{sec:ch05-pure-neumann-problem}
\setcounter{thmcount}{0}
\setcounter{equation}{0}

上一节处理的混合边界条件总含有本质的 Dirichlet 分量. 纯 Neumann(或自然)边界条件
\MBSourceItem{1}
\begin{equation}
\MathBookFormulaID{formula-ch05-pure-neumann-condition}
\label{eq:ch05-pure-neumann-condition}
\frac{\partial u}{\partial\nu}=0\quad\text{于 }\partial\Omega
\end{equation}
(即 $\Gamma=\varnothing$)的情形略有不同, 与一维情形相同(参见练习~\ref{ex:ch00-neumann-not-well-posed}). 特别地, 解只在相差一个加性常数的意义下唯一, 且仅当式~\eqref{eq:ch05-poisson-equation} 的右端 $f$ 满足
\MBSourceItem{2}
\begin{equation}
\MathBookFormulaID{formula-ch05-neumann-compatibility}
\label{eq:ch05-neumann-compatibility}
\int_\Omega f(x)\,dx
=\int_\Omega-\Delta u(x)\,dx
=\int_\Omega\nabla u(x)\cdot\nabla1\,dx
-\int_{\partial\Omega}\frac{\partial u}{\partial\nu}\,ds=0
\end{equation}
时解才可能存在, 这里使用了命题~\ref{prop:ch05-green-identity}. 稍后将看到, 当 $\lambda$ 是简单特征值时, 这种现象对 $-\Delta+\lambda I$ 一类椭圆算子是典型的; 此处 $\lambda=0$. 当前情形适当的变分空间为
\MathBookSourcePage{145}
\MBSourceItem{3}
\begin{equation}
\MathBookFormulaID{formula-ch05-mean-zero-variational-space}
\label{eq:ch05-mean-zero-variational-space}
V=\left\{v\in H^1(\Omega):\int_\Omega v(x)\,dx=0\right\}.
\end{equation}
对任意可积函数 $g$, 定义其平均值 $\overline g$ 为
\MBSourceItem{4}
\begin{equation}
\MathBookFormulaID{formula-ch05-domain-mean}
\label{eq:ch05-domain-mean}
\overline g:=\frac1{\operatorname{meas}(\Omega)}\int_\Omega g(x)\,dx.
\end{equation}
对任意 $v\in H^1(\Omega)$, 有 $v-\overline v\in V$. 使用上一节相同的技术可证明以下命题.

\MBSourceItem{5}
\begin{Proposition}
\label{prop:ch05-pure-neumann-classical-implies-variational}
设 $u\in H^2(\Omega)$ 满足 Poisson 方程~\eqref{eq:ch05-poisson-equation} 和纯 Neumann 边界条件(式~\eqref{eq:ch05-mixed-boundary-conditions} 中 $\Gamma=\varnothing$, 或式~\eqref{eq:ch05-pure-neumann-condition}); 这蕴含 $f\in L^2(\Omega)$ 满足式~\eqref{eq:ch05-neumann-compatibility}. 则 $u-\overline u$ 满足变分表述~\eqref{eq:ch05-poisson-variational-problem}, 其中 $V$ 由式~\eqref{eq:ch05-mean-zero-variational-space} 定义.
\end{Proposition}

其逆命题比上一节的相应结论更复杂.

\MBSourceItem{6}
\begin{Proposition}
\label{prop:ch05-pure-neumann-variational-implies-classical}
设 $f\in L^2(\Omega)$, 且 $u\in H^2(\Omega)$ 满足变分方程~\eqref{eq:ch05-poisson-variational-problem}, 其中 $V$ 由式~\eqref{eq:ch05-mean-zero-variational-space} 定义. 则 $u$ 满足 Poisson 方程~\eqref{eq:ch05-poisson-equation}, 其右端为
\[
\MathBookFormulaID{formula-ch05-mean-corrected-right-hand-side}
\widetilde f(x):=f(x)-\overline f\quad\forall x\in\Omega,
\]
并满足边界条件~\eqref{eq:ch05-pure-neumann-condition}.
\end{Proposition}

\MBSourceItem{7}
\begin{Remark}
\label{rem:ch05-pure-neumann-equivalence}
原问题与变分问题的等价性表述和上一节相似, 但 $V$ 的定义改变了, 并且 $f$ 和 $u$ 上出现约束. 一旦知道 $a(\cdot,\cdot)$ 具有强制性, Riesz 表示定理就保证对任意 $f$ 存在解. 因为 $v\in V$ 蕴含 $\int_\Omega v(x)\,dx=0$, 所以
\[
\MathBookFormulaID{formula-ch05-mean-corrected-functional-equivalence}
\int_\Omega v(x)\widetilde f(x)\,dx
=\int_\Omega v(x)f(x)\,dx\quad\forall v\in V,
\]
故 $f$ 与 $\widetilde f$ 的变分问题相同.
\end{Remark}
\begin{Proof}
上面的注说明无论使用 $f$ 还是 $\widetilde f$, 变分问题都相同. 因此只需假设 $f$ 的平均值为零, 并验证式~\eqref{eq:ch05-poisson-equation} 和式~\eqref{eq:ch05-pure-neumann-condition}. 使用上一节的论证可知
\[
\MathBookFormulaID{formula-ch05-mean-zero-test-equation}
\int_\Omega(-\Delta u-f)v\,dx=0
\quad\forall v\in\widetilde{\mathcal D}(\Omega)
:=\{\phi\in\mathcal D(\Omega):\overline\phi=0\}.
\]

\MathBookSourcePage{146}
容易看出, $\widetilde{\mathcal D}(\Omega)$ 在 $L^2(\Omega)$ 中的闭包是子空间
\[
\MathBookFormulaID{formula-ch05-mean-zero-l2-space}
\widetilde L^2(\Omega):=\{\phi\in L^2(\Omega):\overline\phi=0\}
\]
(参见练习~\ref{exr:ch05-01}). 因而 $-\Delta u=f+\delta$ 在 $L^2(\Omega)$ 中成立, 其中 $\delta=\overline{-\Delta u}$. 再使用命题~\ref{prop:ch05-green-identity}, 得
\[
\MathBookFormulaID{formula-ch05-pure-neumann-boundary-proof}
\begin{aligned}
0&=(f,v)-a(u,v)
=\int_\Omega(-\Delta u-\delta)v\,dx-a(u,v)\\
&=\int_\Omega(-\Delta u)v\,dx-a(u,v)
=-\int_{\partial\Omega}v\frac{\partial u}{\partial\nu}\,ds
\quad\forall v\in V.
\end{aligned}
\]
保持 $v\in V$ 的同时, $v|_{\partial\Omega}$ 显然可以任意选取: 任取 $w\in H^1(\Omega)$, 令 $v:=w-\overline w\phi$, 其中 $\phi\in\mathcal D(\Omega)$ 且 $\overline\phi=1$. 因此式~\eqref{eq:ch05-pure-neumann-condition} 成立. 最后在命题~\ref{prop:ch05-green-identity} 中再取 $v\equiv1$, 得 $\delta=0$.
\end{Proof}
